% Capítulo 3: Sección 3.7 - Trazo de Curvas

Método sistemático y completo para graficar funciones aplicando todos los conceptos del cálculo diferencial: simetrías, asíntotas, crecimiento y concavidad.

La gráfica de una función proporciona una información muy valiosa con respecto al comportamiento de la función. Una de las principales aplicaciones del cálculo radica en proporcionar un procedimiento para el trazo de curvas con una precisión razonable.

Hasta ahora hemos tratado algunos aspectos específicos del trazo de curvas: dominio, contradominio y simetría, en el Capítulo 1; derivadas y tangentes, en el Capítulo 2; y valores extremos, monotonía, concavidad, puntos de inflexión y asíntotas, en este capítulo. Ahora es tiempo de reunir toda esta información para trazar gráficas que revelen las características importantes de las funciones.

\vspace{1.5em}
\section*{Ejercicios}

\noindent \textit{En los Ejercicios 1 a 50 analice las curvas conforme a los apartados A-H.}

\begin{enumerate}
    \begin{minipage}{0.45\textwidth}
    \subimport{ejercicio_01/}{ejercicio.tex}
    \addtocounter{enumi}{1}
    \subimport{ejercicio_03/}{ejercicio.tex}
    \addtocounter{enumi}{1}
    \subimport{ejercicio_05/}{ejercicio.tex}
    \addtocounter{enumi}{1}
    \subimport{ejercicio_07/}{ejercicio.tex}
    \addtocounter{enumi}{1}
    \subimport{ejercicio_09/}{ejercicio.tex}
    \addtocounter{enumi}{1}
    \subimport{ejercicio_11/}{ejercicio.tex}
    \addtocounter{enumi}{1}
    \subimport{ejercicio_13/}{ejercicio.tex}
    \addtocounter{enumi}{1}
    \subimport{ejercicio_15/}{ejercicio.tex}
    \addtocounter{enumi}{1}
    \subimport{ejercicio_17/}{ejercicio.tex}
    \addtocounter{enumi}{1}
    \subimport{ejercicio_19/}{ejercicio.tex}
    \addtocounter{enumi}{1}
    \subimport{ejercicio_21/}{ejercicio.tex}
    \addtocounter{enumi}{1}
    \subimport{ejercicio_23/}{ejercicio.tex}
    \addtocounter{enumi}{1}
    \subimport{ejercicio_25/}{ejercicio.tex}
    \addtocounter{enumi}{1}
    \subimport{ejercicio_27/}{ejercicio.tex}
    \addtocounter{enumi}{1}
    \subimport{ejercicio_29/}{ejercicio.tex}
    \addtocounter{enumi}{1}
    \subimport{ejercicio_31/}{ejercicio.tex}
    \addtocounter{enumi}{1}
    \subimport{ejercicio_33/}{ejercicio.tex}
    \addtocounter{enumi}{1}
    \subimport{ejercicio_35/}{ejercicio.tex}
    \addtocounter{enumi}{1}
    \subimport{ejercicio_37/}{ejercicio.tex}
    \addtocounter{enumi}{1}
    \subimport{ejercicio_39/}{ejercicio.tex}
    \addtocounter{enumi}{1}
    \subimport{ejercicio_41/}{ejercicio.tex}
    \addtocounter{enumi}{1}
    \subimport{ejercicio_43/}{ejercicio.tex}
    \end{minipage}%
    \begin{minipage}{0.45\textwidth}
    \setcounter{enumi}{1}
    \subimport{ejercicio_02/}{ejercicio.tex}
    \addtocounter{enumi}{1}
    \subimport{ejercicio_04/}{ejercicio.tex}
    \addtocounter{enumi}{1}
    \subimport{ejercicio_06/}{ejercicio.tex}
    \addtocounter{enumi}{1}
    \subimport{ejercicio_08/}{ejercicio.tex}
    \addtocounter{enumi}{1}
    \subimport{ejercicio_10/}{ejercicio.tex}
    \addtocounter{enumi}{1}
    \subimport{ejercicio_12/}{ejercicio.tex}
    \addtocounter{enumi}{1}
    \subimport{ejercicio_14/}{ejercicio.tex}
    \addtocounter{enumi}{1}
    \subimport{ejercicio_16/}{ejercicio.tex}
    \addtocounter{enumi}{1}
    \subimport{ejercicio_18/}{ejercicio.tex}
    \addtocounter{enumi}{1}
    \subimport{ejercicio_20/}{ejercicio.tex}
    \addtocounter{enumi}{1}
    \subimport{ejercicio_22/}{ejercicio.tex}
    \addtocounter{enumi}{1}
    \subimport{ejercicio_24/}{ejercicio.tex}
    \addtocounter{enumi}{1}
    \subimport{ejercicio_26/}{ejercicio.tex}
    \addtocounter{enumi}{1}
    \subimport{ejercicio_28/}{ejercicio.tex}
    \addtocounter{enumi}{1}
    \subimport{ejercicio_30/}{ejercicio.tex}
    \addtocounter{enumi}{1}
    \subimport{ejercicio_32/}{ejercicio.tex}
    \addtocounter{enumi}{1}
    \subimport{ejercicio_34/}{ejercicio.tex}
    \addtocounter{enumi}{1}
    \subimport{ejercicio_36/}{ejercicio.tex}
    \addtocounter{enumi}{1}
    \subimport{ejercicio_38/}{ejercicio.tex}
    \addtocounter{enumi}{1}
    \subimport{ejercicio_40/}{ejercicio.tex}
    \addtocounter{enumi}{1}
    \subimport{ejercicio_42/}{ejercicio.tex}
    \end{minipage}
\end{enumerate}

\vspace{1em}
\begin{enumerate}
    \setcounter{enumi}{43}
    \begin{minipage}{0.45\textwidth}
    \subimport{ejercicio_45/}{ejercicio.tex}
    \addtocounter{enumi}{1}
    \subimport{ejercicio_47/}{ejercicio.tex}
    \addtocounter{enumi}{1}
    \subimport{ejercicio_49/}{ejercicio.tex}
    \end{minipage}%
    \begin{minipage}{0.45\textwidth}
    \setcounter{enumi}{43}
    \subimport{ejercicio_44/}{ejercicio.tex}
    \addtocounter{enumi}{1}
    \subimport{ejercicio_46/}{ejercicio.tex}
    \addtocounter{enumi}{1}
    \subimport{ejercicio_48/}{ejercicio.tex}
    \addtocounter{enumi}{1}
    \subimport{ejercicio_50/}{ejercicio.tex}
    \end{minipage}
\end{enumerate}

\vspace{1em}
\noindent \textit{En los Ejercicios 51 a 56 analice cada curva conforme a los apartados A-H. En D encuentre la ecuación de la asíntota inclinada.}

\begin{enumerate}
    \setcounter{enumi}{50}
    \begin{minipage}{0.45\textwidth}
    \subimport{ejercicio_51/}{ejercicio.tex}
    \addtocounter{enumi}{1}
    \subimport{ejercicio_53/}{ejercicio.tex}
    \addtocounter{enumi}{1}
    \subimport{ejercicio_55/}{ejercicio.tex}
    \end{minipage}%
    \begin{minipage}{0.45\textwidth}
    \setcounter{enumi}{51}
    \subimport{ejercicio_52/}{ejercicio.tex}
    \addtocounter{enumi}{1}
    \subimport{ejercicio_54/}{ejercicio.tex}
    \addtocounter{enumi}{1}
    \subimport{ejercicio_56/}{ejercicio.tex}
    \end{minipage}
\end{enumerate}

\vspace{1em}
\begin{enumerate}
    \setcounter{enumi}{56}
    \subimport{ejercicio_57/}{ejercicio.tex}
    \subimport{ejercicio_58/}{ejercicio.tex}
\end{enumerate}
